\subsection{正规基定理}

设 N 是 K 的伽罗瓦扩张，其伽罗瓦群为 \(\Gamma\) 。我们把 \(\Gamma\) 与群代数 \(\mathbf{K}^{(\Gamma)}\) 的典范基等同（III，第446页）；于是 N 可被视为左 \(\mathbf{K}^{(\Gamma)}\) -模（III，第447页，例），从而

\[
u. x = \sum_ {\sigma \in \Gamma} a _ {\sigma} \sigma (x) \quad \text { for } \quad x \in \mathbf {N} \quad \text { and } \quad u = \sum_ {\sigma \in \Gamma} a _ {\sigma} \sigma \quad \text { in } \quad \mathbf {K} ^ {(\Gamma)}.
\]

如果 \(\mathbf{N}\) 在 \(\mathbf{K}\) 上是有限次的，则由戴德金定理（V，第27页，推论2），群 \(\Gamma\) 是有限的，因此我们可以定义元素 \(t = \sum_{\sigma \in \Gamma} \sigma\) （在 \(\mathbf{K}^{(\Gamma)}\) 中）；于是我们有

\[
\operatorname{Tr} _ {N / K} (x) = \sum_ {\sigma \in \Gamma} \sigma (x),
\]

也就是说， \(\operatorname{Tr}_{\mathbf{N} / \mathbf{K}}(x) = t.x\) 对所有 \(x \in \mathbf{N}\) 成立。

让我们定义 \(\Gamma\) 在 \(\mathbf{N}\) 上的一个右作用，即 \(x^{\sigma} = \sigma^{-1}(x)\) 。类似地，我们可以把乘法群 \(\mathbf{N}^{*}\) 视为右 \(\mathbf{Z}^{(\Gamma)}\) -模，其外法则写作 \((x, u) \mapsto x^u\) 。例如，记号 \(x^{2\sigma + 3\tau + \pi}\) （其中 \(\sigma, \tau, \pi\) 是 \(\Gamma\) 的元素）表示乘积 \((x^\sigma)^2 \cdot (x^\tau)^3 \cdot x^\pi\) 。若 \(\mathbf{N}\) 在 \(\mathbf{K}\) 上是有限次的，且 \(t = \sum_{\sigma \in \Gamma} \sigma\) 如上所述，则我们有 \(\mathrm{N}_{\mathrm{N/K}}(x) = \prod_{\sigma \in \Gamma} x^\sigma\) ，即 \(\mathrm{N}_{\mathrm{N/K}}(x) = x^t\) 对所有 \(x \in \mathbf{N}^*\) 成立。

此后假设 N 在 K 上是有限次的。给定 \(x \in N\) ， \(\{x\}\) 构成 \(\mathbf{K}^{(\Gamma)}\) -模 N 的一组基，当且仅当族 \((\sigma(x))_{\sigma \in \Gamma}\) 是 N 在 K 上的一组基。这样的基称为 N 在 K 上的正规基。

定理 6. --- 设 N 是 K 上有限次的伽罗瓦扩张，且令 \(\Gamma\) 为 N 在 K 上的伽罗瓦群。则 N 在 K 上存在一个正规基；换言之， \(\mathbf{K}^{(\Gamma)}\) -模 N 是秩为 1 的自由模。

我们将给出该命题的两个证明。第一个证明使用以下引理，该引理将在第八章（§ 2，第 5 号）中证明。

* 引理 3. --- 设 A 是一个 K-代数， \(\mathbf{M}_1\) 和 \(\mathbf{M}_2\) 是（在 K 上）秩有限的两个 A-模，并假设存在 K 的一个扩域 L，使得模 \(\mathrm{L} \otimes_{\mathrm{K}} \mathrm{M}_1\) 和 \(\mathrm{L} \otimes_{\mathrm{K}} \mathrm{M}_2\) 在环 \(\mathrm{L} \otimes_{\mathrm{K}} \mathrm{A}\) 上是同构的。则 A-模 \(\mathbf{M}_1\) 和 \(\mathbf{M}_2\) 是同构的。

我们将在 \(\mathbf{A} = \mathbf{K}^{(\Gamma)}\) ， \(\mathbf{M}_1 = \mathbf{N}\) ， \(\mathbf{M}_2 = \mathbf{A}_s\) 且 \(\mathbf{L} = \mathbf{N}\) 的情形下应用引理 3。由 V 第64页的推论，存在一个 K-同构 \(\varphi\) （从 \(\mathbf{N} \otimes_{\mathbf{K}} \mathbf{N}\) 到 \(\mathbf{N} \otimes_{\mathbf{K}} \mathbf{K}^{(\Gamma)}\) 上），它把 \(x \otimes y\) 映为 \(\sum x \sigma^{-1}(y) \otimes \sigma\) 。显然 \(\varphi\) 是 \(\mathbf{N} \otimes_{\mathbf{K}} \mathbf{K}^{(\Gamma)}\) -模的同构，因此定理由引理 3 得出。 *

对于第二个证明，我们将使用以下命题：

命题 10. --- 设 \(x \in \mathbb{N}\) ，则 \(\{x\}\) 构成 \(K^{(\Gamma)}\) -模 \(\mathbb{N}\) 的一组基，其充分必要条件是 \(\det(\sigma\tau(x))_{\sigma,\tau \in \Gamma} \neq 0\) 。

由于 \(\mathbf{K}^{(\Gamma)}\) 和 \(\mathbf{N}\) 在 \(\mathbf{K}\) 上有相同的维数，说 \(\{x\}\) 是 \(\mathbf{N}\) 在 \(\mathbf{K}^{(\Gamma)}\) 上的一组基，意味着映射 \(a \mapsto ax\) （从 \(\mathbf{K}^{(\Gamma)}\) 到 \(\mathbf{N}\) 中）是单射。这又意味着映射 \(b \mapsto b(1 \otimes x)\) （从 \(\mathbf{N} \otimes_{\mathbf{K}} \mathbf{K}^{(\Gamma)}\) 到 \(\mathbf{N} \otimes_{\mathbf{K}} \mathbf{N}\) 中）是单射（II，第306页，命题14）。现在存在一个 \(\mathbf{N} \otimes_{\mathbf{K}} \mathbf{K}^{(\Gamma)}\) -模同构（从 \(\mathbf{N} \otimes_{\mathbf{K}} \mathbf{N}\) 到 \(\mathbf{N} \otimes_{\mathbf{K}} \mathbf{K}^{(\Gamma)}\) 上），它把 \(1 \otimes x\) 映为 \(\sum_{\sigma} \sigma^{-1}(x) \otimes \sigma\) 。由此可知， \(\{x\}\) 是 N 在

\(\mathbf{K}^{(\Gamma)}\) 上的一组基，当且仅当对于 N 中元素的每一个非零族 \((n_{\tau})_{\tau\in\mathrm{T}}\) ，我们有 \(\left(\sum n_{\tau}\otimes\tau\right)\left(\sum\sigma^{-1}(x)\otimes\sigma\right)\neq0\) 。但后一个关系意味着存在 \(\sigma\in\Gamma\) 使得 \(\sum_{\tau}n_{\tau}\sigma^{-1}\tau(x)\neq0\) ，由此得出该命题。

\begin{enumerate}
\def\labelenumi{\Alph{enumi})}
\tightlist
\item
  假设 K 是无限的；映射 \(x \mapsto \det(\sigma\tau(x))\) 是 N 到 N 中的一个（在 K 上的）多项式映射（IV，第54页）。通过从 K 到 N 的标量扩张，我们得到向量 N-空间 \(\mathbf{N} \otimes_{\mathbf{K}} \mathbf{N}\) 上的相应映射，并且我们刚刚看到后者不恒为零（因为 \(\mathbf{N} \otimes_{\mathbf{K}} \mathbf{N}\) 在 \(\mathbf{N} \otimes_{\mathbf{K}} \mathbf{K}^{(\Gamma)}\) 上是秩为 1 的自由模）。所以存在 \(x \in \mathbf{N}\) 使得 \(\det(\sigma\tau(x)) \neq 0\) （IV，第18页，定理2）；更一般地，由同一文献我们有：
\end{enumerate}

命题 11. --- 假设 K 是无限的，设 P:N→K 是一个在 K 上非零的多项式映射。存在 x∈N 使得 P(x)≠0 且 \{x\} 是 N 在 K 上的一组基 \(^{(\Gamma)}\) 。

\begin{enumerate}
\def\labelenumi{\Alph{enumi})}
\setcounter{enumi}{1}
\tightlist
\item
  假设 K 是有限的。由命题4（V，第95页） \(^{1}\) ，K 的每一个有限次扩张都有循环伽罗瓦群。因此我们将更一般地考虑群 \(\Gamma\) 是 n 阶循环群的情形；我们用 \(\gamma\) 表示 \(\Gamma\) 的一个生成元。
\end{enumerate}

下述引理是第七章中证明的更一般结果的特殊情形。环 A 或者是有理整数环 Z，或者是域 K 上的多项式环 K{[}X{]} 。

引理 4. --- 设 M 是一个由有限个元素 \(x_{1}, \ldots, x_{h}\) 生成的挠 A-模；则存在 M 的一个元素 x，其零化子（II，第219页）等于 M 的零化子。

在这两种情形下，A 都是一个整环，且 A 的每一个理想都是主理想。当 A = Z（相应地 A = K{[}X{]}）时，我们用 P 表示素数集合（相应地 K{[}X{]} 中不可约首一多项式的集合）。对于 A 的每一个元素 \(a \neq 0\) ，存在 A 的一个可逆元素 u 和一个具有有限支撑的正整数族 \((v_{p}(a))_{p \in \mathcal{P}}\) ，使得 \(a = u \prod_{p \in \mathcal{P}} p^{v_{p}(a)}\) 且 u 和整数

\(v_{p}(a)\) 是唯一确定的（I，第51页和IV，第13页，命题13）。

设 \(\mathfrak{a}_i\) 为 \(x_i\) 的零化子（对于 \(1 \leqslant i \leqslant h\) ）， \(\mathfrak{a}\) 为 \(\mathbf{M}\) 的零化子；设 \(a_1, \ldots, a_h\) ， \(a\) 是 \(\mathbf{A}\) 的非零元素，使得 \(\mathfrak{a}_i = \mathbf{A}a_i\) 且 \(\mathfrak{a} = \mathbf{A}a\) ；由于 \(\mathfrak{a} = \mathfrak{a}_1 \cap \ldots \cap \mathfrak{a}_h\) ，由前述可知

\[
v _ {p} (a) = \sup _ {1 \leqslant i \leqslant h} v _ {p} (a _ {i}) \quad \text { for   all } p \in \mathscr {P}.\tag{10}
\]

我们把 \(a\) 写成形式 \(up_1^{n(1)} \ldots p_r^{n(r)}\) ，其中 \(p_1, \ldots, p_r\) 在 \(\mathcal{P}\) 中互不相同， \(n(1) > 0, \ldots, n(r) > 0\) ，且 \(u\) 是 \(A\) 的可逆元。设 \(j = 1, \ldots, r\) ；由（10），存在一个整数 \(c(j)\) 使得 \(1 \leqslant c(j) \leqslant h\) 且 \(v_{p_j}(a_{c(j)}) = n(j)\) ；存在 \(b_j\) （在 \(A\) 中）满足 \(a_{c(j)} = p_j^{n(j)}b_j\) ，且元素 \(y_j = b_jx_{c(j)}\) 的零化子为理想 \(Ap_j^{n(j)}\) 。

让我们证明 \(y = y_{1} + \cdots + y_{r}\) 的零化子 b 等于 M 的零化子 a。无论如何我们有 \(a \subset b\) ，所以 b 具有形式 \(Ap_{1}^{m(1)} \ldots p_{r}^{m(r)}\) ，其中 \(0 \leqslant m(j) \leqslant n(j)\) ， \(1 \leqslant j \leqslant r\) 。如果我们有 \(a \neq b\) ，那么将存在一个整数 j 使得 \(1 \leqslant j \leqslant r\) 且 \(m(j) < n(j)\) ，因此 \(d_{j} = a/p_{j}\) 将零化 y。现在我们有 \(d_{j}y_{k} = 0\) （对 \(k \neq j\) ），由此我们得到 \(d_{j}y_{j} = 0\) ；但 \(y_{j}\) 的零化子是 \(Ap_{j}^{m(j)}\) ，而 \(d_{j}\) 不是 \(p_{j}^{m(j)}\) 的倍元。因此假设 \(a \neq b\) 是荒谬的。

我们将把引理4应用于以下情形： \(\mathbf{A}\) 是多项式环 \(\mathbf{K}[\mathbf{X}]\) ， \(\mathbf{M}\) 是阿贝尔群 \(\mathbf{N}\) ，其外法则由 \(a \cdot x = \sum_{k=0}^{\infty} c_k \gamma^k(x)\) 给出，其中 \(a = \sum_{k=0}^{\infty} c_k \mathbf{X}^k\) 在 \(\mathbf{K}[\mathbf{X}]\) 中，且 \(x \in \mathbb{N}\) 。设 \(\mathfrak{a}\) 为 \(\mathbf{M}\) 的零化子，则 \(\gamma^n = 1\) ，因此多项式 \(\mathbf{X}^n - 1\) 属于 \(\mathfrak{a}\) 。设 \(\mathbf{F} \in \mathfrak{a}\) ；由IV，第11页，推论，存在元素 \(c_0, c_1, \ldots, c_{n-1}\) （属于 \(\mathbf{K}\) ）和 \(\mathbf{G} \in \mathbf{K}[\mathbf{X}]\) 使得

\[
\mathrm{F} (\mathbf {X}) = c _ {0} + c _ {1} \mathbf {X} + \dots + c _ {n - 1} \mathbf {X} ^ {n - 1} + (\mathbf {X} ^ {n} - 1) \mathrm{G} (\mathbf {X}).\tag{11}
\]

于是我们有 \(c_{0} + c_{1}\gamma + \cdots + c_{n-1}\gamma^{n-1} = 0\) ，它在 \(\operatorname{Hom}_{\mathbf{K}}(\mathbf{N}, \mathbf{N})\) 中成立，并且由于 N 的自同构 1， \(\gamma\) ， \(\gamma^{2}\) ， \ldots, \(\gamma^{n-1}\) 互不相同，戴德金定理（V，第27页，推论2）蕴含 \(c_{0} = c_{1} = \cdots = c_{n-1} = 0\) 。最后，我们有

\[
\mathrm{F} (\mathbf {X}) = \left(\mathbf {X} ^ {n} - 1\right) \mathrm{G} (\mathbf {X}), \quad \text { that   is }, \quad \mathfrak {a} = \left(\mathbf {X} ^ {n} - 1\right) \mathrm{K} [ \mathbf {X} ].
\]

由引理4，存在一个元素 \(x\) （属于 \(\mathbf{N}\) ），它在 \(\mathbf{K}[\mathbf{X}]\) 中的零化子等于 \((\mathbf{X}^n - 1)\mathbf{K}[\mathbf{X}]\) 。由于单项式 1, \(\mathbf{X}, \ldots, \mathbf{X}^{n-1}\) 构成 \(\mathbf{K}[\mathbf{X}]\) 的一个与 \((\mathbf{X}^n - 1)\mathbf{K}[\mathbf{X}]\) 互补的向量子空间的一组基（IV，第11页，推论），元素 \(x, \gamma(x), \ldots, \gamma^{n-1}(x)\) （属于 \(\mathbf{N}\) ）在 \(\mathbf{K}\) 上线性无关。由于 \([\mathbf{N}:\mathbf{K}] = n(\mathbf{V},\mathbf{p}.66,\mathbf{Th}.3)\) ，序列 \((x,\gamma(x),\ldots,\gamma^{n-1}(x))\) 因此是 \(\mathbf{N}\) 在 \(\mathbf{K}\) 上的一个（正规）基。

